A learned robot policy is never only a model. I came to understand that its behavior depends on the semantics and integrity of the surrounding system: whether timestamps align observations with actions, whether state and action fields retain the same meaning from collection to execution, whether datasets pass structural and motion checks, whether a server returns the expected chunks, and whether inference occurs on a meaningful robot state. These dependencies shifted my focus from isolated AI prediction to reliable learning systems for embodied agents.

To gain end-to-end competence across machines and computation, I built my education as a bridge rather than a replacement. I completed a B.Eng. in Mechanical Design, Manufacturing and Automation at Harbin Institute of Technology, Weihai, in June 2025. That training grounded me in mechanism behavior, kinematics, and physical constraints. I am now pursuing a second bachelor's degree in Computer Science and Technology at Harbin Institute of Technology, Shenzhen, expected in June 2027. This expansion lets me reason about both physical capability and how data structures, algorithms, interfaces, and runtime decisions determine machine behavior. Together, they help me trace a learned capability from representation to physical execution.

My research asks how action representation shapes what a vision-language-action model can learn. As a Research Assistant at iLearn, HIT Shenzhen, I am second author of Hermite Action Tokenizer for VLA, an unsubmitted working paper, and I am responsible for the discrete autoregressive tokenization branch. I worked on splitting seven-dimensional action chunks into six-dimensional end-effector trajectories and gripper states, using shared breakpoints to remove redundant endpoints and quantization seams while coding gripper changes separately. I also compared adaptive MLP-based and fixed Hermite segmentation. In an offline codec benchmark averaging four LIBERO subsets with 8-by-7 action chunks, Hermite's Endpoint MSE was approximately 99.7\% lower than FAST and 82.6\% lower than BEAST. I read this result as evidence about codec fidelity under a defined representation and protocol, not as rollout success, inference accuracy, or physical-robot performance. A change in encoding mattered only if its evaluation remained interpretable. That distinction taught me to ask which layer a metric actually measures before treating it as evidence for a learning system. Representation design determines what a model can express, while evaluation determines what conclusions its numbers can support.

Complementary systems work on a Dobot CR5 showed me where algorithmic evaluation stops. I anchored an HDF5 synchronization pipeline to RGB timestamps and aligned robot states, target actions, and gripper feedback through binary search and linear interpolation. Before an episode entered downstream use, I checked schema and count consistency, frame drops, motion discontinuities, workspace bounds, start poses, and slow playback. I then converted synchronized episodes into LeRobot v2.0, preserving seven-dimensional robot-plus-gripper state and action semantics, and adapted OpenPI with training configuration, WebSocket policy serving, and a CR5 inference client. Deployment tests exposed behavior no offline codec metric could certify: non-blocking execution allowed new inference on an intermediate in-motion state and produced back-and-forth corrections, so I used blocking short-horizon ServoP action chunks to preserve state-action consistency. Hermite and CR5 were separate efforts, but together they clarified complementary responsibilities. Algorithmic evaluation can establish how an encoding behaves under a specified benchmark; systems reliability requires clocks, schemas, service contracts, observations, and execution semantics to remain coherent through the actual loop.

Supporting projects reinforced this discipline. For a weld-seam HMI, I integrated C++/Qt controls with Python/PyTorch segmentation inference through QProcess and OpenGL visualization. I preserved boundaries between the desktop interface, inference process, and robot representation instead of collapsing them into one opaque application. During the Xbotics Embodied AI Community Internship, I migrated an ANYmal-C navigation task from Isaac Lab into a MotrixLab NumPy environment. Treating the port as an interface-preservation problem, I refactored reset and step flows and added rollout and reward/termination regression checks. Across both projects, explicit contracts and tests became my tools for cross-component change.

My immediate goal is AI/ML or learning-systems R\&D for embodied agents, building methods that remain dependable when connected to real data and hardware. I want graduate training to deepen my foundations in machine learning, data systems, and scalable software while testing them against sensing and control constraints. After gaining stronger real-system experience, I may consider doctoral study as an option rather than a predetermined step. I seek curricular depth and applied opportunities to turn this perspective into rigorous, transferable engineering.
